% gen_event Architecture Diagram — Event Manager with Multiple Handlers
% Shows the core 1-to-N handler model unique to gen_event.
% Compile with: xelatex gen_event_architecture.tex
\documentclass[border=10pt]{standalone}

% ============================================================
% Packages
% ============================================================
\usepackage{fontspec}
\usepackage{xeCJK}
\setCJKmainfont[BoldFont=STHeiti, ItalicFont=STKaiti]{STSong}
\setCJKsansfont{STHeiti}
\setCJKmonofont{STFangsong}

\usepackage{xcolor}
\usepackage{tikz}
\usepackage{pifont}

\usetikzlibrary{arrows.meta, positioning, shapes.geometric, fit, calc, backgrounds, decorations.pathreplacing}

% ============================================================
% Color Definitions
% ============================================================
\definecolor{erlred}{HTML}{A4070A}
\definecolor{erlblue}{HTML}{2B5EA7}
\definecolor{erlgreen}{HTML}{2E7D32}
\definecolor{erlyellow}{HTML}{F9A825}
\definecolor{erlpurple}{HTML}{6A1B9A}
\definecolor{erlmodule}{HTML}{D84315}
\definecolor{erlinit}{HTML}{00897B}
\definecolor{erlpid}{HTML}{E65100}
\definecolor{erlname}{HTML}{6A1B9A}
\definecolor{erlteal}{HTML}{00695C}
\definecolor{erlhandler}{HTML}{1565C0}
\definecolor{erlevent}{HTML}{E65100}
\definecolor{erlmanager}{HTML}{283593}
\definecolor{darkgray}{HTML}{333333}

% ============================================================
% Document
% ============================================================
\begin{document}

\begin{tikzpicture}[scale=0.72, every node/.style={scale=0.72}]

    % ==================================================================
    % LEFT SIDE: Event Sources / Clients
    % ==================================================================
    \node[draw=erlpurple!60, fill=erlpurple!3, rounded corners=6pt,
          minimum width=7cm, minimum height=12.5cm] (leftframe) at (-5, 0) {};
    \node[font=\small\bfseries\color{erlpurple}, anchor=south] at (leftframe.north)
        {事件源 / Client 进程};

    % --- Top-Left: Manager Lifecycle API ---
    \node[draw=erlgreen, fill=erlgreen!8, rounded corners=5pt,
          minimum width=6.5cm, minimum height=3.8cm] (lifecycle) at (-5, 3.2) {};
    \node[font=\small\bfseries\color{erlgreen!80!black}, anchor=north west]
        at ([xshift=3pt, yshift=-3pt]lifecycle.north west) {Manager 生命周期};
    \node[font=\scriptsize\ttfamily, text width=5.8cm, align=left, anchor=north west]
        at ([xshift=6pt, yshift=-18pt]lifecycle.north west) {
        \textcolor{erlgreen!70!black}{\textbf{启动：}}\\
        gen\_event:start\_link(\textcolor{erlname}{\textbf{EventMgrRef}})\\
        gen\_event:start\_link()\\[4pt]
        \textcolor{erlgreen!70!black}{\textbf{Handler 管理：}}\\
        gen\_event:add\_handler(\textcolor{erlname}{\textbf{Mgr}}, \textcolor{erlmodule}{\textbf{Handler}}, Args)\\
        gen\_event:add\_sup\_handler(\textcolor{erlname}{\textbf{Mgr}}, \textcolor{erlmodule}{\textbf{Handler}}, Args)\\
        gen\_event:delete\_handler(\textcolor{erlname}{\textbf{Mgr}}, \textcolor{erlmodule}{\textbf{Handler}}, Args)\\
        gen\_event:swap\_handler(\textcolor{erlname}{\textbf{Mgr}}, \{Old,A1\}, \{New,A2\})
    };

    % --- Bottom-Left: Event Sending API ---
    \node[draw=erlevent!80, fill=erlevent!8, rounded corners=5pt,
          minimum width=6.5cm, minimum height=4.2cm] (eventsrc) at (-5, -1.8) {};
    \node[font=\small\bfseries\color{erlevent!80!black}, anchor=north west]
        at ([xshift=3pt, yshift=-3pt]eventsrc.north west) {事件发送 API};
    \node[font=\scriptsize\ttfamily, text width=5.8cm, align=left, anchor=north west]
        at ([xshift=6pt, yshift=-18pt]eventsrc.north west) {
        \textcolor{erlevent!80!black}{\textbf{异步广播（不等待）：}}\\
        gen\_event:\textcolor{erlred}{\textbf{notify}}(\textcolor{erlname}{\textbf{Mgr}}, Event)\\[4pt]
        \textcolor{erlevent!80!black}{\textbf{同步广播（等待所有Handler处理完）：}}\\
        gen\_event:\textcolor{erlblue}{\textbf{sync\_notify}}(\textcolor{erlname}{\textbf{Mgr}}, Event)\\[4pt]
        \textcolor{erlevent!80!black}{\textbf{定向调用（指定Handler，同步）：}}\\
        gen\_event:\textcolor{erlpurple}{\textbf{call}}(\textcolor{erlname}{\textbf{Mgr}}, \textcolor{erlmodule}{\textbf{Handler}}, Req)\\[4pt]
        \textcolor{erlevent!80!black}{\textbf{停止：}}\\
        gen\_event:stop(\textcolor{erlname}{\textbf{Mgr}})
    };

    % --- Annotation below left frame ---
    \node[font=\tiny\color{gray}, anchor=north west, text width=6.2cm, align=left]
        at ([xshift=4pt, yshift=-8pt]eventsrc.south west) {
        \ding{72} 任何进程都可以是事件源\\
        \ding{72} notify 是异步的（cast 语义）\\
        \ding{72} sync\_notify 是同步的（等所有 handler 返回）\\
        \ding{72} call 是定向同步调用（只到一个 handler）
    };

    % ==================================================================
    % RIGHT SIDE: Event Manager Process (single process, multiple handlers)
    % ==================================================================
    \node[draw=erlmanager!60, fill=erlmanager!3, rounded corners=6pt,
          minimum width=8.5cm, minimum height=12.5cm] (rightframe) at (6.5, 0) {};
    \node[font=\small\bfseries\color{erlmanager}, anchor=south] at (rightframe.north)
        {Event Manager 进程（单一进程）};
    \node[font=\small\ttfamily\bfseries\itshape\color{black}, anchor=north west]
        at ([xshift=3pt, yshift=-1pt]rightframe.north west) {gen\_event.erl};

    % --- Manager Core ---
    \node[draw=erlblue, fill=erlblue!8, rounded corners=5pt,
          minimum width=8cm, minimum height=2.8cm] (mgrcore) at (6.5, 4.2) {};
    \node[font=\small\bfseries\color{erlblue}, anchor=north west]
        at ([xshift=3pt, yshift=-3pt]mgrcore.north west) {Manager 核心（消息循环）};
    \node[font=\scriptsize, text width=7.2cm, align=left, anchor=north west]
        at ([xshift=6pt, yshift=-18pt]mgrcore.north west) {
        \textcolor{erlblue!80!black}{\ding{72} \textbf{receive 循环：}}\\
        {\tiny 接收 notify/call/info 消息 $\rightarrow$ 遍历 Handler 列表}\\[3pt]
        \textcolor{erlblue!80!black}{\ding{72} \textbf{内部状态：}}\\
        {\tiny \texttt{\#state\{name, handlers, hibernate\_after\}}}\\
        {\tiny \texttt{handlers} = [\texttt{\#handler\{module, id, state\}}, ...]}
    };

    % --- Handler List (the key visual: multiple handlers in one process) ---
    \node[draw=erlteal!70, fill=erlteal!5, rounded corners=5pt,
          minimum width=8cm, minimum height=7.2cm] (hlist) at (6.5, -1.2) {};
    \node[font=\small\bfseries\color{erlteal!80!black}, anchor=north west]
        at ([xshift=3pt, yshift=-3pt]hlist.north west) {Handler 列表（同一进程内，顺序执行）};

    % Handler A
    \node[draw=erlhandler!80, fill=erlhandler!10, rounded corners=4pt,
          minimum width=7.2cm, minimum height=1.6cm] (ha) at (6.5, 0.8) {};
    \node[font=\scriptsize\bfseries\color{erlhandler}, anchor=north west]
        at ([xshift=3pt, yshift=-2pt]ha.north west)
        {Handler A: \textcolor{erlmodule}{\texttt{log\_handler}}};
    \node[font=\tiny\ttfamily, text width=6.4cm, align=left, anchor=north west]
        at ([xshift=6pt, yshift=-14pt]ha.north west) {
        State\_A = \#\{file => Fd, level => info\}\\
        回调: init/1, handle\_event/2, handle\_call/2, terminate/2
    };

    % Handler B
    \node[draw=erlhandler!80, fill=erlhandler!10, rounded corners=4pt,
          minimum width=7.2cm, minimum height=1.6cm] (hb) at (6.5, -1.2) {};
    \node[font=\scriptsize\bfseries\color{erlhandler}, anchor=north west]
        at ([xshift=3pt, yshift=-2pt]hb.north west)
        {Handler B: \textcolor{erlmodule}{\texttt{stats\_handler}}};
    \node[font=\tiny\ttfamily, text width=6.4cm, align=left, anchor=north west]
        at ([xshift=6pt, yshift=-14pt]hb.north west) {
        State\_B = \#\{count => 42, ets\_tab => Tab\}\\
        回调: init/1, handle\_event/2, handle\_call/2, terminate/2
    };

    % Handler C
    \node[draw=erlhandler!80, fill=erlhandler!10, rounded corners=4pt,
          minimum width=7.2cm, minimum height=1.6cm] (hc) at (6.5, -3.2) {};
    \node[font=\scriptsize\bfseries\color{erlhandler}, anchor=north west]
        at ([xshift=3pt, yshift=-2pt]hc.north west)
        {Handler C: \textcolor{erlmodule}{\texttt{alarm\_handler}}};
    \node[font=\tiny\ttfamily, text width=6.4cm, align=left, anchor=north west]
        at ([xshift=6pt, yshift=-14pt]hc.north west) {
        State\_C = \#\{alarms => [], sup\_pid => Pid\}\\
        回调: init/1, handle\_event/2, handle\_call/2, terminate/2
    };

    % --- Annotation: each handler has independent state ---
    \node[font=\tiny\color{erlteal!80!black}, anchor=north west, text width=7.2cm, align=left]
        at ([xshift=6pt, yshift=2pt]hlist.south west) {
        \ding{72} 每个 Handler 有\textbf{独立的 State}，互不干扰\\
        \ding{72} 所有 Handler 在\textbf{同一进程}中顺序执行（非并行）\\
        \ding{72} 一个 Handler 崩溃 $\rightarrow$ 仅移除该 Handler，不影响其他
    };

    % ==================================================================
    % ARROWS
    % ==================================================================

    % Arrow ①: start_link (lifecycle -> manager core)
    \draw[-{Stealth[length=6pt]}, ultra thick, erlgreen!80!black, line width=1.6pt]
        ([yshift=8mm]lifecycle.east) --
        node[above, font=\scriptsize, align=center, text=erlgreen!80!black]
        {\textcolor{black}{\ding{172}} \texttt{start\_link}\\[1pt]
         {\tiny 创建 Manager 进程}}
        ([yshift=8mm]mgrcore.west);

    % Arrow ②: add_handler (lifecycle -> handler list)
    \draw[-{Stealth[length=6pt]}, thick, erlteal!80!black, line width=1.4pt]
        ([yshift=-8mm]lifecycle.east) --
        node[below, font=\scriptsize, align=center, text=erlteal!80!black, yshift=-2pt]
        {\textcolor{black}{\ding{173}} \texttt{add\_handler}\\[1pt]
         {\tiny 安装 Handler 到列表}}
        ([yshift=22mm]hlist.west);

    % Arrow ③: notify — FAN-OUT (the key visual!)
    % Main arrow from eventsrc to manager
    \draw[-{Stealth[length=6pt]}, ultra thick, erlred!80!black, line width=1.8pt]
        ([yshift=6mm]eventsrc.east) --
        node[above, font=\scriptsize, align=center, text=erlred!80!black]
        {\textcolor{black}{\ding{174}} \texttt{notify / sync\_notify}\\[1pt]
         {\tiny 广播 Event 到所有 Handler}}
        ([yshift=6mm]hlist.west);

    % Fan-out arrows inside the handler list (from left edge to each handler)
    \coordinate (fanpoint) at ([xshift=4mm]hlist.west);
    \draw[-{Stealth[length=4pt]}, thick, erlred!60, line width=1pt]
        ([xshift=4mm, yshift=6mm]hlist.west) -- ([xshift=-2mm]ha.west);
    \draw[-{Stealth[length=4pt]}, thick, erlred!60, line width=1pt]
        ([xshift=4mm, yshift=6mm]hlist.west) -- ([xshift=-2mm]hb.west);
    \draw[-{Stealth[length=4pt]}, thick, erlred!60, line width=1pt]
        ([xshift=4mm, yshift=6mm]hlist.west) -- ([xshift=-2mm]hc.west);

    % Arrow ④: call — directed to ONE handler
    \draw[-{Stealth[length=5pt]}, thick, erlpurple, line width=1.2pt, dashed]
        ([yshift=-10mm]eventsrc.east) --
        node[below, font=\scriptsize, align=center, text=erlpurple, yshift=-2pt]
        {\textcolor{black}{\ding{175}} \texttt{call}\\[1pt]
         {\tiny 定向到指定 Handler}}
        ([xshift=-2mm]hb.west);

    % Arrow ⑤: Reply from call
    \draw[-{Stealth[length=5pt]}, thick, erlpid!80, line width=1.2pt, dashed]
        ([xshift=-2mm, yshift=-4mm]hb.west) --
        node[below, font=\scriptsize, text=erlpid!80, yshift=-1pt]
        {\textcolor{black}{\ding{176}} Reply}
        ([yshift=-16mm]eventsrc.east);

    % ==================================================================
    % COMPARISON NOTE at bottom
    % ==================================================================
    \node[draw=erlyellow!80, fill=erlyellow!10, rounded corners=4pt,
          minimum width=15cm, minimum height=1.8cm, anchor=north]
        (comparison) at (0.75, -6.8) {};
    \node[font=\small\bfseries\color{erlyellow!80!black}, anchor=north west]
        at ([xshift=4pt, yshift=-3pt]comparison.north west) {与 gen\_server 的核心区别};
    \node[font=\tiny, text width=14cm, align=left, anchor=north west]
        at ([xshift=6pt, yshift=-16pt]comparison.north west) {
        \textbf{gen\_server}：一个进程 = 一个回调模块 = 一份 State \hfill
        \textbf{gen\_event}：一个进程 = \textcolor{erlred}{\textbf{N 个回调模块}} = \textcolor{erlred}{\textbf{N 份独立 State}}\\[2pt]
        gen\_server 的 \texttt{call} 由唯一的 \texttt{handle\_call/3} 处理 \hfill
        gen\_event 的 \texttt{notify} 会\textcolor{erlred}{\textbf{扇出}}到所有 Handler 的 \texttt{handle\_event/2}\\[2pt]
        gen\_server 崩溃 = 整个服务不可用 \hfill
        gen\_event 单个 Handler 崩溃 = 仅移除该 Handler，Manager 继续运行
    };

\end{tikzpicture}

\end{document}
